{{دیوان
| عنوان = شکوہ
| ماخذ = ویکی ماخذ
| ماخذ_لنک = https://ur.wikisource.org/wiki/%D8%B4%DA%A9%D9%88%DB%81
| مقابلہ = اقبال سائبر لائبریری (اقبال اکادمی پاکستان)
| مقابلہ_لنک = https://www.iqbalcyberlibrary.net
}}
<poem>
کیوں زیاں کار بنوں، سُود فراموش رہوں
فکرِ فردا نہ کروں، محوِ غمِ دوش رھوں

نالے بُلبُل کے سُنوں اور ہمہ تن گوش رہوں
ہم نوا میں بھی کوئی گُل ہوں کہ خاموش رہوں

جُرأت آموز میری تابِ سخن ہے مجھ کو
شکوہ اللہ سے، خاکمِ بدہن ہے مجھ کو

ہے بجا شیوۂ تسلیم میں مشہور ہیں ہم
قصّۂ درد سُناتے ہیں کہ مجبور ہیں ہم

سازِ خاموش ہیں، فریاد سے معمور ہیں ہم
نالہ آتا ہے اگر لب پہ تو معذور ہیں ہم

اے خدا! شکوۂ اربابِ وفا بھی سُن لے
خوگرِ حمد سے تھوڑا سا گِله بھی سُن لے

تھی توموجود ازل سے ہی تری ذاتِ قدیم
پھُول تھا زیبِ چمن پر نہ پریشاں تھی شمیم

شرطِ اِنصاف ہے اَے صاحبِ الطافِ عمیم
بُوئے گُل پھیلتی کس طرح جو ہوتی نہ نسیم

ہم کو جمعیّتِ خاطر یہ پریشانی تھی
ورنہ اُمّت تیرے محبوبؐ کی دیوانی تھی؟

ہم سے پہلے تھا عجب تیرے جہاں کا منظر
کہیں مسجود تھے پتّھر، کہیں معبود شجر

خُوگرِ پیکرِ محسوس تھی اِنساں کی نظر
مانتا پھر کوئی اَن دیکھے خدا کو کیونکر

تجھ کو معلوم ہے لیتا تھا کوئی نام ترا؟
قوّتِ بازوئے مسلم نے کیا کام ترا

بس رہے تھے یہیں سلجوق بھی ،تُورانی بھی
اہلِ چیں چین میں، ایران میں ساسانی بھی

اِسی معمورے میں آباد تھے یونانی بھی
اِسی دُنیا میں یہودی بھی تھے،نصرانی بھی

پر ترے نام پہ تلوار اُٹھائی کس نے
بات جو بگڑی ہوئی تھی، وہ بنائی کس نے

تھے ہمیں ایک ترے معرکہ آراؤں میں
خشکیوں میں کبھی لڑتے ،کبھی دریاؤں میں

دِیں اذانیں کبھی یورپ کے کلیساؤں میں
کبھی افریقہ کے تپتے ہوئے صحراؤں میں

شان آنکھوں میں نہ جچتی تھی جہاں داروں کی
کلمہ پڑھتے تھے ہم چھاؤں میں تلواروں کی

ہم جو جیتے تھے تو جنگوں کی میصبت کے لیے
اور مرتے تھے ترے نام کی عظمت کے لیے

تھی نہ کچھ تیغ زنی اپنی حکومت کے لیے
سربکف پھرتے تھے کیا دہر میں دولت کے لیے؟

قوم اپنی جو زر و مالِ جہاں پر مرتی
بُت فروشی کے عِوض بُت شکنی کیوں کرتی؟

ٹل نہ سکتے تھے اگر جنگ میں اڑ جاتے تھے
پاؤں شیروں کے بھی میداں سے اُکھڑ جاتے تھے

تجھ سے سر کش ہوا کوئی تو بگڑ جاتے تھے
تیغ کیا چیز ہے ہم توپ سے لڑ جاتے تھے

نقش توحید کاہر دِل پہ بٹھایا ہم نے
زیرِ خنجر بھی یہ پیغام سُنایا ہم نے

تُو ہی کہہ دے کہ اُکھاڑا درِ خیبر کس نے
شہر قیصر کا جوتھا،اُسکو کیا سَر کس نے

توڑے مخلوق خداوندوں کے پیکر کس نے
کاٹ کر رکھ دیے کُفّار کے لشکر کس نے

کس نے ٹھنڈا کیا آتشکدۂ ایراں کو؟
کس نے پھر زندہ کیا تذکرۂ یزداں کو؟

کون سی قوم فقط تیری طلبگار ہوئی
اور تیرے لیے زحمت کشِ پیکار ہوئی

کِس کی شمشیر جہاں گیر،جہاں دار ہوئی
کس کی تکبیر سے دنیا تری بیدار ہوئی

کس کی ہیبت سے صنم سہمے ہوئے رہتے تھے
مُنہ کے بَل گر کے "ھُوَ اللہُ اَحَد"کہتے تھے

آگیا عین لڑائی میں اگر وقتِ نماز
قبلہ رُو ہو کے زمیں بوس ہوئی قومِ حجاز

ایک ہی صف میں کھڑے ہو گئے محمود و ایاز
نہ کوئی بندہ رہا اور نہ هي کوئی بندہ نواز

بندہ و صاحب و محتاج و غنی ایک ہوئے
تیری سرکار میں پہنچے تو سبھی ایک ہوئے

محفلِ کون و مکاں میں سحر و شام پھرے
مئے توحید کو لے کر صفتِ جام پھرے

کوہ میں، دشت میں لے کر ترا پیغام پھرے
اور معلوم ہے تجھ کو،کبھی ناکام پھرے

دشت تو دشت ہیں ،دریا بھی نہ چھوڑے ہم نے
بحرِ ظلمات میں دوڑا دیے گھوڑے ہم نے

صفحۂ دہر سے باطل کو مِٹایا ہم نے
نوعِ انساں کو غلامی سے چھڑایا ہم نے

تیرے کعبے کو جبینوں سے بسایا ہم نے
تیرے قرآن کو سینوں سے لگایا ہم نے

پھر بھی ہم سے یہ گِلا ہے کہ وفادار نہیں
ہم وفادار نہیں،تُو بھی تو دِلدار نہیں!

اُمتیں اور بھی ہیں، اُن میں گناہگار بھی ہیں
عجز والے بھی ہیں، مست مئے پندار بھی ہیں

ان میں کاہل بھی ہیں،غافل بھی ہیں،ہشیار بھی ہیں
سینکڑوں ہیں کہ ترے نام سے بیزار بھی ہیں

رحمتیں ہیں تری اغیار کے کاشنوں پر
برق گرتی ہے تو بیچارے مسلمانوں پر

بُت صنم خانوں میں کہتے ہیں، مسلمان گئے
ہے خوشی ان کو کہ کعبے کے نگہبان گئے

منزلِ دہر سے اونٹوں کے حُدی خوان گئے
اپنی بغلوں میں دبائے ہوئے قرآن گئے

خندہ زن ہے کُفر،احساس تجھے ہے کے نہیں
اپنی توحید کا کچھ پاس تجھے ہے کہ نہیں

یہ شکایت نہیں، ہیں اُن کے خزانے معمور
نہیں محفل میں جنھیں بات بھی کرنے کا شعور

قہر تو یہ ہے کہ کافر کو ملیں حُور و قصور
اور بیچارے مسلماں کو فقط وعدۂ حور

اب وہ الطاف نہیں، ہم پہ عنایات نہیں
بات یہ کیا ہے کہ پہلی سی مدارات نہیں

کیوں مسلمانوں میں ہے دولتِ دنیا نایاب
تیری قدرت تو ہے وہ جس کی نہ حد ہے نہ حساب

تُو جو چاہے تو اُٹھے سینۂ صحرا سے حباب
رہروِ دشت ہو سیلی زدۂ موجِ سراب

طعنِ اغیار ہے، رسوائی ہے،ناداری ہے
کیا تیرے نام پہ مرنے کا عوض خواری ہے؟

بنی اغیار کی اب چاہنے والی دنیا
رہ گئی اپنے لیے ایک خیالی دنیا

ہم تو رخصت ہوئے اوروں نے سنبھالی دنیا
پھر نہ کہنا ہوئی توحید سے خالی دنیا

ہم تو جیتے ہیں کہ دنیا میں تیرا نام رہے
کہیں ممکن ہے کہ ساقی نہ رہے، جام رہے؟

تیری محفل بھی گئی، چاہنے والے بھی گئے
شب کی آہیں بھی گئیں، صبح کے نالے بھی گئے

دل تجھے دے بھی گئے، اپنی صِلا لے بھی گئے
آ کے بیٹھے بھی نہ تھے اور نکالے بھی گئے

آئے عشّاق، گئے وعدۂ فردا لے کر
اب اُنہیں ڈھونڈ چراغِ رُخِ زیبا لے کر

دردِ لیلیٰ بھی وہی، قیس کا پہلو بھی وہی
نجد کے دشت و جبل میں رَمِ آہو بھہ وہی

عشق کا دل بھی وہی، حُسن کا جادو بھی وہی
اُمّتِ احمدِ مرسل بھی وہی، تُو بھی وہی

پھر یہ آزردگیِ غیر سبب کا معنی
اپنے شیداؤں پہ یہ چشمِ غضب کیا معنی

تجھ کو چھوڑا کہ رسولِ عربی کو چھوڑا؟
بُت گری پیشہ کیا، بت شکنی کو چھوڑا؟

عشق کو، عشق کی آشفتہ سری کو چھوڑا؟
رسمِ سلمانؓو اویس قرنیؓ کو چھوڑا؟

آگ تکبیر کی سینوں میں دبی رکھتے ہیں
زندگی مثلِ بلال حبشی رکھتے ہیں

عشق کی خیر وہ پہلے سی ادا بھی نہ سہی
جادہ پیمائیِ تسلیم و رضا بھی نہ سہی

مُضطِرب دل صفت قبلہ نما بھی نہ سہی
اور پابندیِ آئینِ وفا بھی نہ سہی

کبھی ہم سے، کبھی غیروں سے شناسائی ہے
بات کہنے کی نہیں،تُو بھی تو ہرجائی ہے!

سرِ فاراں پہ کِیا دین کو کامل تو نے
اک اشارے میں ہزاروں کے لیے دل تو نے

آتش اندوز کیا عشق کا حاصل تُو نے
پُھونک دی گرمیِ رُخسار سے محفل تُو نے

آج کیوں سینے ہمارے شرر آباد نہیں
ہم وہی سوختہ ساماں ہیں، تجھے یاد نہیں؟

وادی نجد میں وہ شورِ سلاسل نہ رہا
قیس دیوانۂ نظّارۂ محمل نہ رہا

حوصلے وہ نہ رہے، ہم نہ رہے، دل نہ رہا
گھر یہ اُجڑا ہے کہ تُو رونقِ محفل نہ رہا

اے خوش آں روز کہ آئی و بصد ناز آئی
بے حجابانہ سُوئے محفلِ ما باز آئی

بادہ کش غیر ہیں گلشن میں لبِ جُو بیٹھے
سُنتے ہیں جام بکف نغمۂ کُو کُو بیٹے

دور ہنگامۂ گُلزار سے یک سُو بیٹھے
تیرے دیوانے بھی ہیں مُنتظر "ھُو" بیٹھے

اپنے پروانوں کی پھر ذوقِ خُود افروزی دے
برقِ دیرینہ کو فرمانِ جگر سوزی دے

قوم آوارہ عناں تاب ہے پھر سُوئے حجاز
لے اُڑا بُلبلِ بے پر کو مذاقِ پرواز

مضطرب باغ کے ہر غنچے میں ہے، بُوئے نیاز
تُو ذرا چھیڑ تو دے، تشنۂ مضراب ہے ساز

نغمے بے تاب ہیں تاروں سے نکلنے کے لیے
طُور مضطر ہے اُسی آگ میں جلنے کے لیے

مُشکلیں اُمّتِ مرحوم کی آساں کر دے
مُورِ بے مایہ کو ہمدوشِ سلیماں کر دے

جنس نایابِ محبّت کو پھر ارزاں کر دے
ہند کے دَیر نشینوں کو مسلماں کر دے

جُوئے خوں می چکدا از حسرتِ دیرینہ ما
می تپد نالۂ بہ نشتر کدۂ سینۂ ما

بُوئے گُل لے گئی بیرون چمن رازِ چمن
کیا قیامت ہے کہ خود پھُول ہیں غمّازِ چمن!

عہدِ گُل ختم ہوا، ٹوٹ گیا سازِ چمن
اُڑ گئے ڈالیوں سے زمزمہ پرداز چمن

ایک بُلبل ہے کہ محو ترنّم ہے اب تک
اس کے سینے میں ہے نغموں کا تلاطم اب تک

قُمریاں شاخِ صنوبر سے گریزاں بھی ہوئیں
پتّیاں پُھول کی جھڑ جھڑ کے پریشاں بھی ہوئیں

وہ پرانی روشیں باغ کی ویراں بھی ہوئیں
ڈالیاں پیرہنِ برگ سے عُریاں بھی ہوئیں

قیدِ موسم سے طبیعت رہی آزاد اس کی
کاش گُلشن میں سمجھتا کوئی فریاد اس کی!

لُطف مرنے میں ہے باقی نہ مزا جینے میں
کچھ مزا ہے تو یہی خونِ جگر پینے میں

کتنے بیتاب ہیں جوہر مرے آئینے میں
کس قدر جلوے تڑپتے ہیں مرے سینے میں

اس گُلستاں میں مگر دیکھنے والے ہی نہیں
داغ جو سینے میں رکھتا ہوں، وہ لالے ہی نہیں

چاک اس بُلبلِ تنہا کی نوا سے دل ہوں
جاگنے والے اسی بانگِ درا سے دل ہوں

یعنی پھر زندہ نئے عہدِ وفا سے دل ہوں
پھر اِسی بادۂ دیرینہ کے پیاسے دل ہوں

عجمی خُم ہے تو کیا،مے تو حجازی ہے مری
نغمہ ہندی ہے تو کیا،لَے تو حجازی ہے مری
</poem>
